\chapter{Concluzii}

Proiectul Mini PC folosind modulul LattePanda Mu fezabilitatea proiectării unui PCB complex și a unui sistem integrat performan cu un grad ridicat de personalizare, capabil să ruleze sisteme de operare și aplicații moderne. Arhitectura hardware oferă o flexibilitate ridicată atât din punct de vedere al modificării PCB-ului existent cu alte funcționalități dar și prin posibilitatea de schimbare a modulului cu unul mai performant. 

\textbf{În viitor}, următoarele modificări trebuie aduse platformei pentru a putea fi produs în masă:
\begin{itemize}
  \item Optimizarea carcasei pentru turnare prin injecție, acest lucru oferind un timp foarte scurt de producție pentru elementele de plastic, printarea 3D a acestor durând 6-8 ore pe o imprimantă standard, timp redus la cateva minute pentru turnarea prin injecție. De asemenea piesele rezultate sunt mai rezistente;
  \item Modificarea prinderilor PCB-ului secundar, pentru montare mai rapidă pep liniile de producție.
\end{itemize}

Platforma actuală reprezintă un punct de plecare solid, atât pentru producție în masă cât și pentru dezvoltarea altor proiecte pe baza sa. 